\subsection{Transaction Cost Analysis (TCA)}

\begin{definition} \hlt{TCA}\\
Transaction cost analysis compares realised execution prices with relevant benchmarks and explains the gap by market context, order design, broker behaviour, and trading style.
\end{definition}

\begin{taxonomy} \hlt{Transaction Cost Component Segmentation}
\begin{enumerate}[label=\roman*.]
\setlength{\itemsep}{0pt}
\item \hlt{explicit costs}: commissions, exchange fees, taxes, clearing, settlement, and borrow charges;
\item \hlt{implicit costs}: bid-ask spread, market impact, delay cost, and adverse selection relative to a benchmark;
\item \hlt{missed trade opportunity}: value lost when order is not filled, or filled too slowly to capture intended price.
\end{enumerate}
Explicit costs are accounting entries. Implicit and missed-opportunity costs require a counterfactual benchmark.
\end{taxonomy}

\begin{model} \hlt{Signed Benchmark Cost}\\
Let $d=+1$ for a buy and $d=-1$ for a sale, $q$ be filled quantity, $\bar{p}$ the average fill price, and $b$ the benchmark price. The signed benchmark cost is $C=dq(\bar{p}-b)$.\\
The sign convention makes high buy prices and low sale prices costly.
\end{model}

\begin{model} \hlt{Implementation Shortfall with Missed Quantity}\\
For parent size $Q$, filled quantity $q$, average fill $\bar{p}$, decision benchmark $p_0$, and later evaluation price  $p_T$,\\
$IS=dq(\bar{p}-p_0)+d(Q-q)(p_T-p_0)+\text{explicit costs}$.\\
The first term is realised execution cost. The second term is missed trade opportunity, and it is meaningful only when the unfilled order was economically feasible.
\end{model}

\begin{principle} \hlt{Marginal Aggressiveness Rule}\\
Trade more aggressively when the marginal missed-opportunity cost saved exceeds the marginal transaction cost paid. Trade less aggressively when the marginal transaction cost saved exceeds the additional missed-opportunity cost. Average cost is the wrong control variable for the last slice.
\end{principle}

\begin{taxonomy} \hlt{Benchmark Price Segmentation}
\begin{enumerate}[label=\roman*.]
\setlength{\itemsep}{0pt}
\item \hlt{contemporaneous midpoint}: clean for effective spread, weak for timing skill;
\item \hlt{post-trade midpoint}: measures realised spread and adverse selection;
\item \hlt{decision price}: natural for implementation shortfall and paper-portfolio comparison;
\item \hlt{VWAP}: compares to average traded market price, but can hide own impact and split-order bias;
\item \hlt{open or close}: useful for sponsors with coarse timestamps, but noisy and gameable.
\end{enumerate}
No benchmark is neutral; it defines which cost is visible.
\end{taxonomy}

\begin{definition} \hlt{VWAP Benchmark}\\
For trades $i$ in the benchmark window, $\mathrm{VWAP}=\frac{\sum_i q_i p_i}{\sum_i q_i}=\sum_i w_i p_i$, where $w_i=q_i/\sum_j q_j$. It is an average market execution price, not a no-impact counterfactual.
\end{definition}

\begin{taxonomy} \hlt{Performance Measurement Segmentation}
\begin{enumerate}[label=\roman*.]
\setlength{\itemsep}{0pt}
\item \hlt{absolute}: total return or all-in execution price without a comparison asset;
\item \hlt{relative}: result versus benchmark, peer, arrival price, market volume, or factor model;
\item \hlt{holding-period}: investor-level return after cash flows and distributions;
\item \hlt{money-weighted}: IRR, useful when capital flows are controlled by the investor;
\item \hlt{trade-weighted}: fill-quality measures, useful when the desk controls execution.
\end{enumerate}
Relative measures more useful as most price movement is not caused by trader or manager being evaluated.
\end{taxonomy}

\begin{model} \hlt{Risk-Adjusted Excess Return}\\
For portfolio/strategy return $r^p_t$, benchmark return $r^m_t$, a one-factor evaluation is $r^p_t-r^f_t=\alpha+\beta(r^m_t-r^f_t)+u_t$.\\
Here $\alpha$ is realised selection value after market exposure, not proof of skill. Microstructure TCA uses the same logic at shorter horizons: remove common price movement before attributing residual cost to execution.
\end{model}

\begin{taxonomy} \hlt{Benchmark Failure Modes}
\begin{enumerate}[label=\roman*.]
\setlength{\itemsep}{0pt}
\item \hlt{wrong exposure}: benchmark lacks the risk actually carried;
\item \hlt{wrong clock}: benchmark window does not match the decision or execution horizon;
\item \hlt{own-impact contamination}: the order helps set the benchmark it is judged against;
\item \hlt{gaming}: agents can change behaviour to improve the measured score;
\item \hlt{selection}: only favourable orders, venues, or managers enter the sample.
\end{enumerate}
\end{taxonomy}

\begin{principle} \hlt{Pre-, Live-, and Post-Trade Consistency}\\
The same indicator should have a pre-trade expectation, a live monitored value, and a post-trade assessment. A bid-ask-spread indicator, for example, can move from expected average, to updated live likelihood, to realised influence on performance.
\end{principle}

\begin{taxonomy} \hlt{Performance Attribution Segmentation}
\begin{enumerate}[label=\roman*.]
\setlength{\itemsep}{0pt}
\item \hlt{immediate price}: the reference available when the order decision is made;
\item \hlt{planned price}: the price expected from the risk-control trajectory;
\item \hlt{unexpected market move}: exogenous price variation during execution;
\item \hlt{trading efficiency}: quality of liquidity-seeking tactics relative to the plan;
\item \hlt{market impact}: price movement caused by the order itself.
\end{enumerate}
This segmentation is useful only when the benchmark, horizon, and side are explicit.
\end{taxonomy}

\begin{model} \hlt{Selection and Timing Decomposition}\\
Let $r^p_t$ be realised return, $r^m_t$ market return, and $\beta_t$ the exposure chosen during the period.\\
A useful attribution separates: $r^p_t = r^m_t + (r^p_t-\beta_t r^m_t)+(\beta_t-1)r^m_t$.\\
The three terms are market return, selection residual, and timing payoff from changing exposure. In execution, the analogue is market move, trading residual, and urgency or participation timing.
\end{model}

\begin{caveat} \hlt{TCA Uses Wider Samples Than Post-Trade Review}\\
Post-trade review explains one execution in detail. TCA usually aggregates months of comparable trades to identify contexts where an algorithm, broker, or parameter set is systematically strong or weak.
\end{caveat}

\begin{taxonomy} \hlt{TCA Estimator Quality Segmentation}
\begin{enumerate}[label=\roman*.]
\setlength{\itemsep}{0pt}
\item \hlt{data}: daily benchmarks are cheap; quote methods need synchronised trades and quotes;
\item \hlt{benchmark}: closer references reduce noise but shift timing attribution;
\item \hlt{inference}: splitting, momentum, informed flow, gaming, and small samples bias costs.
\end{enumerate}
\end{taxonomy}

\begin{taxonomy} \hlt{TCA Context Segmentation}
\begin{enumerate}[label=\roman*.]
\setlength{\itemsep}{0pt}
\item \hlt{liquidity}: spread, queue size, turnover, trade size, free float, and trading rate;
\item \hlt{risk}: volatility, jumps, intraday trend, and index or sector co-movement;
\item \hlt{order design}: size, duration, participation rate, benchmark, and limit constraints;
\item \hlt{market state}: news, auctions, rebalances, corporate events, and venue outages.
\end{enumerate}
\end{taxonomy}

\begin{taxonomy} \hlt{Pre-Trade Cost Prediction Variables}
\begin{enumerate}[label=\roman*.]
\setlength{\itemsep}{0pt}
\item \hlt{order variables}: parent size, participation target, limit price, urgency, and placement;
\item \hlt{current market variables}: spread, displayed depth, volatility, recent returns, and signed volume;
\item \hlt{general market variables}: ADV, market capitalisation, risk regime, sector liquidity, and event calendar;
\item \hlt{venue variables}: hidden fill rates, fee tier, queue loss rate, and routing latency.
\end{enumerate}
Prediction is a conditioning problem: yesterday's cost matters only after adjusting for today's liquidity state.
\end{taxonomy}

\begin{practitionernote} \hlt{Broker and Algo Selection}\\
The practical output of TCA is a routing table: for a given market context and trading objective, send the next order to the broker or algorithm whose historical performance is best after adjusting for risk and liquidity.
\end{practitionernote}

\begin{principle} \hlt{Best Execution as Auditability}\\
Best execution is not the cheapest commission or the fastest fill in isolation. It is the best auditable outcome given the client's objective, limit, urgency, available venues, and information risk.
\end{principle}

\begin{taxonomy} \hlt{Retail Execution Quality Segmentation}
\begin{enumerate}[label=\roman*.]
\setlength{\itemsep}{0pt}
\item \hlt{price}: execution relative to NBBO, midpoint, or effective spread;
\item \hlt{speed}: time to fill, cancel, or reject;
\item \hlt{certainty}: fill probability and partial-fill risk;
\item \hlt{commission}: explicit brokerage charge net of rebates;
\item \hlt{adverse selection}: post-fill price movement and opportunity cost.
\end{enumerate}
Retail traders observe commissions easily; execution price quality requires quote-synchronised data.
\end{taxonomy}

\begin{taxonomy} \hlt{Broker-Conflict Segmentation}
\begin{enumerate}[label=\roman*.]
\setlength{\itemsep}{0pt}
\item \hlt{routing inducements}: rebates, payments for order flow, reciprocal flow, and affiliate venues;
\item \hlt{information misuse}: front-running, inappropriate exposure, and prearranged trades;
\item \hlt{allocation}: assigning favourable fills to preferred accounts;
\item \hlt{activity incentives}: commissions, churning, soft dollars, and commission recapture;
\item \hlt{principal conflict}: internalisation, dual trading, and dealer inventory.
\end{enumerate}
TCA and audit trails convert these conflicts from stories into measurable routing and fill-quality checks.
\end{taxonomy}

\begin{caveat} \hlt{Gaming the Measure}\\
Any benchmark that brokers can influence can be gamed. Delaying a trade, matching the benchmark mechanically, splitting fills around measurement windows, or selecting only favourable venues can improve reported TCA while worsening the client's economics.
\end{caveat}

\begin{caveat} \hlt{Smooth Returns Are Not Safer Returns}\\
Illiquid valuations, stale marks, and discretion in marking can reduce reported volatility and raise apparent risk-adjusted performance. The risk has not disappeared; it has moved into valuation lag and jump risk.
\end{caveat}

\begin{process} \hlt{Block Motive Audit}
\begin{enumerate}[label=\roman*.]
\setlength{\itemsep}{0pt}
\item identify economic reason for trade: funding, diversification, hedging, benchmark change, or speculation;
\item test whether the requested size and timing are consistent with that reason;
\item compare the initiator's alternatives, constraints, and past behaviour;
\item price the trade only after separating utilitarian urgency from possible information.
\end{enumerate}
The audit is not moral judgement; it is adverse-selection control.
\end{process}

\begin{caveat} \hlt{Block Trade Reporting Trade-Off}\\
Fast public reporting improves transparency and protects the regular market, but it can hurt dealers that still need to unwind facilitated inventory. Delayed reporting protects facilitation, but gives informed block traders more room to exploit stale prices.
\end{caveat}
